\documentclass[../main.tex]{subfiles}
\ifSubfilesClassLoaded{\setcounter{chapter}{4}}{}
\begin{document}

\chapter{Encoding Kategorikal, Scaling, dan Persiapan Matriks Fitur}

\begin{subcpmk}
  \item \textbf{Sub-CPMK-P3:} Encoding dan scaling; menyusun matriks \texttt{X} dan vektor \texttt{y} tanpa kebocoran informasi yang kasar (RPS Mg.\ 4).
\end{subcpmk}

\SesiRPS{4}{P3}{$\sim$170 menit}{CPMK-P2}

\section{Kemampuan akhir sesi}

Membedakan fitur nominal vs ordinal; menerapkan one-hot atau label encoding; menggunakan \texttt{StandardScaler} atau \texttt{MinMaxScaler}; menjelaskan risiko \textit{leakage} jika scaler difit pada gabungan latih+uji.

\section{Teori}

\begin{konsep}
Model seperti regresi dan SVM umumnya membutuhkan fitur numerik ter-skala. Kategorikal perlu dikodekan. \textbf{Data leakage} terjadi jika statistik dari seluruh data (termasuk uji) dipakai untuk transform sebelum split---pada minggu ini kita siapkan pola yang akan distandarkan di minggu 6+ dengan \texttt{Pipeline} dan split yang benar.
\end{konsep}

\begin{catatan}
Contoh Titanic: \texttt{fit\_transform} hanya pada \texttt{X\_train}; \texttt{transform} pada \texttt{X\_test}. Jangan menggabungkan latih+uji lalu baru memutuskan imputer/scaler.
\end{catatan}

\section{Sarana}

\begin{itemize}[leftmargin=*, nosep]
  \item \texttt{pandas}, \texttt{seaborn}
  \item \texttt{sklearn.preprocessing}, \texttt{sklearn.pipeline}
  \item \texttt{sklearn.compose.ColumnTransformer}
\end{itemize}

\section{Prosedur kerja}

\begin{enumerate}[leftmargin=*]
  \item Pilih subset kolom prediktor dan satu target (jika sudah ada).
  \item Pisahkan \texttt{X} dan \texttt{y}; untuk eksplorasi encoding, dokumentasikan di Markdown.
  \item Terapkan \texttt{OneHotEncoder} atau \texttt{OrdinalEncoder} sesuai konteks.
  \item Skala numerik; gabungkan menjadi \texttt{X\_prepared} (array atau DataFrame).
\end{enumerate}

\section{Contoh kode}

\begin{lstlisting}[caption=Titanic: pisah fitur/target dan subset kolom]
import seaborn as sns

df = sns.load_dataset("titanic")
df = df.drop(columns=["deck"], errors="ignore")

target_col = "survived"
drop_extra = ["class", "who", "adult_male", "embark_town", "alive", "alone"]
X = df.drop(columns=[target_col] + drop_extra, errors="ignore")
y = df[target_col]
\end{lstlisting}

\begin{lstlisting}[caption=Split dulu; prapemroses hanya fit di data latih]
from sklearn.model_selection import train_test_split

X_train, X_test, y_train, y_test = train_test_split(
    X, y, test_size=0.2, random_state=42, stratify=y
)
\end{lstlisting}

\begin{lstlisting}[caption=ColumnTransformer + imputer (polar anti-leakage)]
from sklearn.compose import ColumnTransformer
from sklearn.preprocessing import OneHotEncoder, StandardScaler
from sklearn.impute import SimpleImputer
from sklearn.pipeline import Pipeline

num_features = ["pclass", "age", "sibsp", "parch", "fare"]
cat_features = ["sex", "embarked"]

numeric_pipe = Pipeline([
    ("imp", SimpleImputer(strategy="median")),
    ("scale", StandardScaler()),
])
categorical_pipe = Pipeline([
    ("imp", SimpleImputer(strategy="most_frequent")),
    ("oh", OneHotEncoder(handle_unknown="ignore")),
])

preprocess = ColumnTransformer(
    transformers=[
        ("num", numeric_pipe, num_features),
        ("cat", categorical_pipe, cat_features),
    ]
)

X_train_prep = preprocess.fit_transform(X_train)
X_test_prep = preprocess.transform(X_test)
print("train", X_train_prep.shape, "test", X_test_prep.shape)
\end{lstlisting}

\begin{lstlisting}[caption=Alternatif cepat: pandas get\_dummies (tanpa pipeline)]
import pandas as pd

X_num = X_train[num_features].copy()
X_cat = X_train[cat_features].copy()
X_dum = pd.get_dummies(X_cat, drop_first=True)
X_prep_df = pd.concat([X_num.reset_index(drop=True), X_dum], axis=1)
print(X_prep_df.shape)
\end{lstlisting}

\section{Referensi}

\begin{itemize}[leftmargin=*]
  \item sklearn preprocessing. \url{https://scikit-learn.org/stable/modules/preprocessing.html}
\end{itemize}

\begin{asesmen}
\textbf{Formatif:} narasi encoding/scaling; dimensi \texttt{X\_prep} konsisten; tidak ada target di fitur.
\end{asesmen}

\begin{rangkuman}
Minggu 4 menyiapkan representasi numerik untuk pemodelan dengan kesadaran leakage.

\textbf{Kata kunci:} one-hot, scaler, ColumnTransformer, SimpleImputer
\end{rangkuman}

\end{document}
